\section{Euler 方法：从局部斜率到全局误差}
\label{sec:08-Numerical-Analysis/06-ODE/01}

考虑一个你连微分方程都不需要用就能猜出答案的问题：
\[
y'(t)=-y(t),\qquad y(0)=1.
\]
解是 \(y(t)=e^{-t}\)。假设我们不知道这个公式——或者假设 \(f(t,y)\) 复杂到根本写不出解析解——只能用数值方法一步一步往前推。

已知 \(y(0)=1\)。当 \(t\) 从 0 走到 \(h\)，函数值变了多少？导数 \(y'(0)=-1\) 告诉我们：在起点处，\(y\) 正在以速率 1 向下衰减。走一小步 \(h\)，如果速率保持不变，函数值大约变成
\[
y_1=y_0+h\cdot y'(0)=1+h\cdot(-1)=1-h.
\]
取 \(h=0.5\)，得到 \(y_1=0.5\)。真解 \(e^{-0.5}\approx 0.6065\)——差了一成多。

再把 \((0.5,0.5)\) 当作新的起点，导数 \(y'(0.5)\approx -0.5\)，再往前半步：
\[
y_2=0.5+0.5\cdot(-0.5)=0.25.
\]
真解 \(e^{-1}\approx 0.3679\)。误差从 0.107 变成了 0.118。继续这样走十步，数值解越来越离谱，而 0.5 这个步长看起来并不是一个天文数字。

这就是显式 Euler 方法的最简形式。已知 \((t_n,y_n)\)，下一个点由
\[
y_{n+1}=y_n+h f(t_n,y_n)
\]
给出。用一句话说：它拿当前位置的切线当作下一小段轨迹。每一步的直觉都是对的——在这一点上，往这个方向走，确实暂时沿着解的曲线。但累积起来就出了问题。

\subsection{局部看了没问题，全局为什么只有一阶}

把精确解 \(y(t)\) 在 \(t_n\) 附近做 Taylor 展开：
\[
y(t_n+h)
=y(t_n)+h f(t_n,y(t_n))
+\frac{h^2}{2}y''(\xi_n),
\]
其中 \(\xi_n\) 是 \(t_n\) 和 \(t_n+h\) 之间的某个点。Euler 步恰好取了前两项，扔掉了 \(y''\) 项。如果这一步恰好从精确解 \(y(t_n)\) 出发，Euler 的局部截断误差——单步产生的误差——是 \(O(h^2)\)。二阶看起来不错。

但在固定时间区间 \([t_0,T]\) 上要走
\[
N\approx\frac{T-t_0}{h}
\]
步。每一步都有 \(O(h^2)\) 的新误差注入，总共约 \(N\cdot O(h^2)=O(h)\) 的\"新误差产量\"。更关键的是，旧误差不会停在原地：每一步的起点已经偏离了精确解，导数也就估错了位置，下一个 Euler 步在错误的起点上再往前走——误差会传播。

记 \(e_n=y(t_n)-y_n\) 为第 \(n\) 步时的全局误差。把精确解的 Taylor 展开减去 Euler 递推公式：
\[
e_{n+1}
=e_n+h\bigl[f(t_n,y(t_n))-f(t_n,y_n)\bigr]
+\frac{h^2}{2}y''(\xi_n).
\]
方括号里的差是导数在精确值和近似值处的差异。若 \(f\) 对 \(y\) 满足 Lipschitz 条件（存在常数 \(L\) 使得 \(|f(t,y_1)-f(t,y_2)|\le L|y_1-y_2|\)），且解足够光滑使得 \(|y''|\le 2C\)，则误差递推可以放缩为
\[
|e_{n+1}|
\le(1+hL)|e_n|+Ch^2.
\]
这个递推公式的信息量很大：第一项 \((1+hL)|e_n|\) 是\"上一轮旧误差被动力系统放大\"的部分——\(L\) 越大，误差被放得越狠；第二项 \(Ch^2\) 是\"这一轮新注入的误差\"。

现在从 \(e_0=0\) 出发，反复展开这个不等式：
\[
|e_n|
\le Ch^2\sum_{k=0}^{n-1}(1+hL)^k
=\frac{Ch}{L}\bigl[(1+hL)^n-1\bigr].
\]
利用 \(1+hL\le e^{hL}\)，再用 \(nh\le T-t_0\)：
\[
|e_n|
\le\frac{Ch}{L}\bigl(e^{L(t_n-t_0)}-1\bigr).
\]
在固定区间上，括号里的指数项是有界常数，因此
\[
\max_{t_n\le T}|e_n|=O(h).
\]
局部二阶缺陷积累成全局一阶误差。把局部截断阶和全局阶混为一谈——比如声称\"Euler 的误差是 \(O(h^2)\)\"——是 ODE 方法中最容易发生也最致命的概念混淆。两步的量级差了一个 \(h\)，而 \(h\) 恰好在分母上决定了步数。

\subsubsection{指数衰减的显式验证}

回到 \(y'=-y,\;y(0)=1\)。显式 Euler 从递推 \(y_{n+1}=(1-h)y_n\) 给出
\[
y_n=(1-h)^n,\qquad t_n=nh.
\]
精确解 \(e^{-t_n}\) 与 \((1-h)^n\) 的差，在固定 \(T=nh\) 下随 \(h\) 减半而近似减半——一阶收敛，和上面用 Grönwall 推出来的结论吻合。

但如果 \(h>2\)，事情就完全变了：\(|1-h|>1\)，数值解 \(|y_n|\) 随 \(n\) 增长，而真实解 \(e^{-t}\) 随 \(t\) 衰减。方法突然从\"不太精确\"跳到了\"定性上完全错误\"。这不是截断阶能解释的现象——它属于稳定性的范畴。\hyperref[sec:08-Numerical-Analysis/06-ODE/03]{``稳定区域与刚性方程''}会把这种\"步长大了就会爆炸\"的临界条件变成可计算的判据。

\subsection{一杯咖啡的教训}

一杯热咖啡放在室温中。设温差 \(e=T_{\text{咖啡}}-T_{\text{室温}}\)，牛顿冷却律说温差的变化率正比于温差本身：
\[
e'=-\frac{k}{C}e,
\]
其中 \(k\) 是有效换热强度，\(C\) 是热容量。物理规律保证：温差平滑衰减到零，温度始终单调趋近室温，不会忽冷忽热。

显式 Euler 把这个连续规律离散成：
\[
e_{n+1}
=\left(1-\frac{hk}{C}\right)e_n.
\]
这个离散系统保持单调衰减的充要条件是
\[
\left|1-\frac{hk}{C}\right|<1,
\qquad\text{即}\qquad
0<h<\frac{2C}{k}.
\]
如果 \(hk/C\) 落在 \(1\) 和 \(2\) 之间，数值温差会正负交替：这一秒比室温热，下一秒比室温冷，再下一秒又热——真实的咖啡没有经历任何这样的震荡，是数值方法自己编出来的。如果 \(hk/C>2\)，这种虚构的震荡还会越荡越大，最终让数值温差炸掉。错误不出自热传导，也不出自咖啡，只出自步长——一个和咖啡毫无关系的计算参数，在模拟中获得了扭曲物理的能力。

这个教训比咖啡本身重要得多。一个冷库温控系统用同一个模型做仿真，\(h\) 就变成了控制器读取温度和更新动作的时间间隔。软件里把 \(h\) 设得太大，会在仿真中制造真实冷库并不存在的振荡；如果你在仿真中看到振荡就去调大控制增益来镇压它，你就是在针对一个数值伪影优化真实设备的控制策略。步长在这里从\"画图时选多密\"升级成了一个真正的工程参数：它决定了控制器与物理系统之间的信息采样频率，而不只是曲线是否光滑。

\subsection{后向 Euler：用未来的斜率压住发散}

显式 Euler 的问题是拿着过去的斜率预报未来。如果反过来，拿新点上的斜率来校正自己呢？
\[
y_{n+1}=y_n+h f(t_{n+1},y_{n+1}).
\]
\(y_{n+1}\) 同时出现在等号两边——每步需要解方程（线性的或非线性的），因此称为隐式方法。

对 \(y'=-y\)，方程直接变成 \(y_{n+1}=y_n-h y_{n+1}\)，解得
\[
y_{n+1}=\frac{y_n}{1+h}.
\]
无论 \(h\) 取多大，\(y_{n+1}\) 的绝对值都严格小于 \(|y_n|\) 并且符号不变——物理上单调衰减的定性行为被完整保留了下来。这就是隐式方法在刚性方程中的价值：一步更贵（需要解方程），但换来了对快速衰减模态的强阻尼。

后向 Euler 和显式 Euler 都是一阶方法——全局误差都是 \(O(h)\)。两者的精度级别相同，稳定区域却完全不同。\"隐式\"不等于\"自动精确\"：每步的非线性 Newton 迭代、Jacobian 矩阵的近似质量和线性求解器的容差，都会进入总误差。精度没有免费午餐，但稳定性有时候花点计算量可以买回来。

\subsection{三个容易搞混的概念}

ODE 数值方法的理论大厦主要靠三根柱子撑着：

\begin{itemize}
\tightlist
\item
  一致性：把精确解代入方法的递推公式时，局部缺陷随 \(h\to0\) 而消失。这是\"公式在极限下是对的\"的最低保证。
\item
  稳定性：局部扰动（包括舍入误差和截断误差）不会在离散推进中被不合理地放大，至少在被允许的步长范围内。
\item
  收敛性：固定时间区间上，数值解随 \(h\to0\) 逼近精确解。
\end{itemize}

三者之间的关系可以这样理解：一致性高意味着方法在无限小步长下准确；稳定性好意味着把步长从无限小放宽到有限时不会爆炸；两者兼备才得到收敛——有限步长下的实际逼近。一个一致但不稳定的方法随着步数增加，误差可能指数发散；一个稳定但步长过大的方法，可能给出稳定但不准确的轨迹。

\subsection{实际 ODE 计算中误差的来源不止一种}

书中推导误差界时常假设只有时间离散是误差来源。真实求解器面前列着一整排误差：

\begin{itemize}
\tightlist
\item
  初值和参数本身的测量或建模误差——输入就不精确；
\item
  时间离散误差——就是我们一直在讨论的局部截断和全局累积；
\item
  每步浮点舍入误差——步数越多，舍入累积越多；
\item
  隐式方法内层求解的迭代误差——Newton 迭代也不可能跑到数学收敛；
\item
  事件定位误差——想停在 \(g(t,y)=0\) 的瞬间，实际上只能停在附近；
\item
  模型本身与物理现实的差距——ODE 是对现实的近似，数值方法是近似的近似。
\end{itemize}

把步长减到极小只能压低第二项。步长过小反而会让第三步的舍入累积吃掉所有收益，计算时间也跟着步数的 \(1/h\) 增长。各项误差之间有一个 trade-off 地形，步长选择需要找到地形中的谷底，而不是一味往\"更小\"的极限里钻。

\subsection{如何自己验证阶数}

如果你实现了一个 ODE 方法想验证它的阶数是否正确，做法很直接：找一个有解析解的充分光滑问题，在同一个终止时间 \(T\) 上用三组步长 \(h, h/2, h/4\) 分别计算。若方法已进入渐近区，全局误差 \(E(h)\) 应大致满足
\[
\frac{E(h)}{E(h/2)}\approx 2^p,
\]
\(p\) 是方法的全局阶数。Euler 应得比值约 2。测试时要保证所有步长都在稳定区域内——否则你测到的可能是不稳定发散的速度而不是截断阶，那是个完全不同的物理量。
